% GENERATED FILE. Edit canonical lessons, then run npm run generate:books.
% canonical-source-hash: fnv1a64:f4a28839c08d64e8
% canonical-lessons: KA-C328-darode, KA-C328-docu, KA-C328-esagu, KA-C328-parariyagu, KA-C328-jadu

\chapter{A Robbery at Night}
\label{ch:ka-a-robbery-at-night}
\hlchaptermodality{328}

\begin{chapteropening}
\textbf{By the end of this chapter:} \emph{I can say a robbery, to loot, to commit, to abscond and a track.}

Complete the last lesson of chapter 328: I can say a robbery, to loot, to commit, to abscond and a track.
\end{chapteropening}

\section[\texorpdfstring{\kn{ದರೋಡೆ} (darōḍe)}{ದರೋಡೆ (darōḍe)}]{\kn{ದರೋಡೆ} (darōḍe) — a robbery}

\label{lesson:KA-C328-darode}



\begin{quote}
Before the new one: say the Kannada for abruptly, then the Kannada for to be startled.
\end{quote}

\subsection*{You'll want to know: \kn{ದರೋಡೆ}}
\textbf{\kn{ದರೋಡೆ}} — \emph{darōḍe} — \textquotedblleft{}a robbery\textquotedblright{}.

\begin{quote}
\kn{ನಿನ್ನೆ} \kn{ರಾತ್ರಿ} \kn{ಬ್ಯಾಂಕಿನಲ್ಲಿ} \kn{ದರೋಡೆ} \kn{ಆಯಿತು}.
\end{quote}

\begin{quote}
\emph{ninne rātri byāṅkinalli darōḍe āyitu.}
\end{quote}

\begin{quote}
There was a robbery at the bank last night.
\end{quote}

\kn{ದರೋಡೆ} is robbery by force, often by a gang.

\begin{cousinweb}
\textbf{\kn{ದರೋಡೆ}} (\emph{darōḍe}) is usually traced to Marathi \emph{daroḍā}, \textquotedblleft{}a gang robbery\textquotedblright{}; the earlier history of the word is not settled.
\end{cousinweb}

\begin{grammarlens}[title={What you've built}]
One more word for telling what happened.
\end{grammarlens}

\subsection*{Your turn}
\begin{itemize}
\raggedright
  \item \emph{Say it:} \emph{darōḍe}
  \item \emph{Say it:} \emph{darōḍe}, once more
  \item \emph{Say it:} \emph{ninne rātri byāṅkinalli darōḍe āyitu.} — There was a robbery at the bank last night.
  \item \emph{Recall:} say the Kannada for abruptly, then the Kannada for to be startled, then say \emph{darōḍe} again
\end{itemize}

\begin{tcolorbox}[breakable,colback=teal!4,colframe=teal!35!black,title={Before you move on}]
What does \kn{ದರೋಡೆ} mean? (\textquotedblleft{}A robbery\textquotedblright{}.) Say it once more.
\end{tcolorbox}
\section[\texorpdfstring{\kn{ದೋಚು} (dōcu)}{ದೋಚು (dōcu)}]{\kn{ದೋಚು} (dōcu) — to loot; to rob}

\label{lesson:KA-C328-docu}



\begin{quote}
Before the new one: say the Kannada for to be startled, then the Kannada for a robbery.
\end{quote}

\subsection*{You'll want to know: \kn{ದೋಚು}}
\textbf{\kn{ದೋಚು}} — \emph{dōcu} — \textquotedblleft{}to loot; to rob\textquotedblright{}.

\begin{quote}
\kn{ಕಳ್ಳರು} \kn{ಅಂಗಡಿಯನ್ನು} \kn{ದೋಚಿದರು}.
\end{quote}

\begin{quote}
\emph{kaḷḷaru aṅgaḍiyannu dōcidaru.}
\end{quote}

\begin{quote}
The robbers looted the shop.
\end{quote}

Past (I …): \kn{ದೋಚಿದೆ}, \emph{dōcide}.

The place or person robbed takes -\kn{ಅನ್ನು}: \kn{ಮನೆಯನ್ನು} \kn{ದೋಚು}.

\begin{cousinweb}
\textbf{\kn{ದೋಚು}} (\emph{dōcu}) is native Dravidian: DEDR 3541 lists Kannada \emph{dōcu}, \textquotedblleft{}to plunder, rob\textquotedblright{}, beside Telugu \emph{dō̃cu}, \textquotedblleft{}to plunder, pillage\textquotedblright{}.
\end{cousinweb}

\begin{grammarlens}[title={What you've built}]
One more word for telling what happened.
\end{grammarlens}

\subsection*{Your turn}
\begin{itemize}
\raggedright
  \item \emph{Say it:} \emph{dōcu}
  \item \emph{Say it:} \emph{dōcu}, once more
  \item \emph{Say it:} \emph{kaḷḷaru aṅgaḍiyannu dōcidaru.} — The robbers looted the shop.
  \item \emph{Recall:} say the Kannada for to be startled, then the Kannada for a robbery, then say \emph{dōcu} again
\end{itemize}

\begin{tcolorbox}[breakable,colback=teal!4,colframe=teal!35!black,title={Before you move on}]
What does \kn{ದೋಚು} mean? (\textquotedblleft{}To loot; to rob\textquotedblright{}.) Say it once more.
\end{tcolorbox}
\section[\texorpdfstring{\kn{ಎಸಗು} (esagu)}{ಎಸಗು (esagu)}]{\kn{ಎಸಗು} (esagu) — to commit (a crime)}

\label{lesson:KA-C328-esagu}



\begin{quote}
Before the new one: say the Kannada for a robbery, then the Kannada for to loot.
\end{quote}

\subsection*{You'll want to know: \kn{ಎಸಗು}}
\textbf{\kn{ಎಸಗು}} — \emph{esagu} — \textquotedblleft{}to commit (a crime)\textquotedblright{}.

\begin{quote}
\kn{ಅವನು} \kn{ಹಣಕ್ಕಾಗಿ} \kn{ಆ} \kn{ಅಪರಾಧ} \kn{ಎಸಗಿದನು}.
\end{quote}

\begin{quote}
\emph{avanu haṇakkāgi ā aparādha esagidanu.}
\end{quote}

\begin{quote}
He committed that crime for money.
\end{quote}

Past (I …): \kn{ಎಸಗಿದೆ}, \emph{esagide}.

A formal verb for wrong deeds: \kn{ಅಪರಾಧ} \kn{ಎಸಗು}, \kn{ತಪ್ಪು} \kn{ಎಸಗು}.

\begin{cousinweb}
\textbf{\kn{ಎಸಗು}} (\emph{esagu}) is native Dravidian: DEDR 471 lists Kannada \emph{esagu}, \textquotedblleft{}to engage in, undertake, do, perform\textquotedblright{}.
\end{cousinweb}

\begin{grammarlens}[title={What you've built}]
One more word for telling what happened.
\end{grammarlens}

\subsection*{Your turn}
\begin{itemize}
\raggedright
  \item \emph{Say it:} \emph{esagu}
  \item \emph{Say it:} \emph{esagu}, once more
  \item \emph{Say it:} \emph{avanu haṇakkāgi ā aparādha esagidanu.} — He committed that crime for money.
  \item \emph{Recall:} say the Kannada for a robbery, then the Kannada for to loot, then say \emph{esagu} again
\end{itemize}

\begin{tcolorbox}[breakable,colback=teal!4,colframe=teal!35!black,title={Before you move on}]
What does \kn{ಎಸಗು} mean? (\textquotedblleft{}To commit (a crime)\textquotedblright{}.) Say it once more.
\end{tcolorbox}
\section[\texorpdfstring{\kn{ಪರಾರಿಯಾಗು} (parāriyāgu)}{ಪರಾರಿಯಾಗು (parāriyāgu)}]{\kn{ಪರಾರಿಯಾಗು} (parāriyāgu) — to abscond; to flee}

\label{lesson:KA-C328-parariyagu}



\begin{quote}
Before the new one: say the Kannada for to loot, then the Kannada for to commit.
\end{quote}

\subsection*{You'll want to know: \kn{ಪರಾರಿಯಾಗು}}
\textbf{\kn{ಪರಾರಿಯಾಗು}} — \emph{parāriyāgu} — \textquotedblleft{}to abscond; to flee\textquotedblright{}.

\begin{quote}
\kn{ದರೋಡೆಯ} \kn{ನಂತರ} \kn{ಕಳ್ಳರು} \kn{ಕಾರಿನಲ್ಲಿ} \kn{ಪರಾರಿಯಾದರು}.
\end{quote}

\begin{quote}
\emph{darōḍeya nantara kaḷḷaru kārinalli parāriyādaru.}
\end{quote}

\begin{quote}
After the robbery the robbers fled in a car.
\end{quote}

Past (I …): \kn{ಪರಾರಿಯಾದೆ}, \emph{parāriyāde}.

A news word for an escape from the police: \kn{ಆರೋಪಿ} \kn{ಪರಾರಿಯಾದನು}.

\begin{cousinweb}
\textbf{\kn{ಪರಾರಿಯಾಗು}} (\emph{parāriyāgu}) is \kn{ಪರಾರಿ}, \textquotedblleft{}on the run\textquotedblright{}, with \emph{āgu} \textquotedblleft{}to become\textquotedblright{} (DEDR 333). \kn{ಪರಾರಿ} is usually traced to Hindustani \emph{farārī}, \textquotedblleft{}a fugitive\textquotedblright{}, from Arabic \emph{firār}, \textquotedblleft{}flight\textquotedblright{}.
\end{cousinweb}

\begin{grammarlens}[title={What you've built}]
One more word for telling what happened.
\end{grammarlens}

\subsection*{Your turn}
\begin{itemize}
\raggedright
  \item \emph{Say it:} \emph{parāriyāgu}
  \item \emph{Say it:} \emph{parāriyāgu}, once more
  \item \emph{Say it:} \emph{darōḍeya nantara kaḷḷaru kārinalli parāriyādaru.} — After the robbery the robbers fled in a car.
  \item \emph{Recall:} say the Kannada for to loot, then the Kannada for to commit, then say \emph{parāriyāgu} again
\end{itemize}

\begin{tcolorbox}[breakable,colback=teal!4,colframe=teal!35!black,title={Before you move on}]
What does \kn{ಪರಾರಿಯಾಗು} mean? (\textquotedblleft{}To abscond; to flee\textquotedblright{}.) Say it once more.
\end{tcolorbox}
\section[\texorpdfstring{\kn{ಜಾಡು} (jāḍu)}{ಜಾಡು (jāḍu)}]{\kn{ಜಾಡು} (jāḍu) — a track; a trail}

\label{lesson:KA-C328-jadu}



\begin{quote}
Before the new one: say the Kannada for to commit, then the Kannada for to abscond.
\end{quote}

\subsection*{You'll want to know: \kn{ಜಾಡು}}
\textbf{\kn{ಜಾಡು}} — \emph{jāḍu} — \textquotedblleft{}a track; a trail\textquotedblright{}.

\begin{quote}
\kn{ಪೊಲೀಸರು} \kn{ಕಾರಿನ} \kn{ಜಾಡು} \kn{ಹಿಡಿದು} \kn{ಹೋದರು}.
\end{quote}

\begin{quote}
\emph{polīsaru kārina jāḍu hiḍidu hōdaru.}
\end{quote}

\begin{quote}
The police followed the car's trail.
\end{quote}

\kn{ಜಾಡು} \kn{ಹಿಡಿ}, to follow a trail. Not to be confused with \kn{ಜಾಡಿ}, a jar.

\begin{cousinweb}
\textbf{\kn{ಜಾಡು}} (\emph{jāḍu}) is native Dravidian: DEDR 2442 lists Kannada \emph{jāḍu}, \textquotedblleft{}the mark of a footstep or of a cart wheel, a track, a trace\textquotedblright{}.
\end{cousinweb}

\begin{grammarlens}[title={What you've built}]
One more word for telling what happened.
\end{grammarlens}

\subsection*{Your turn}
\begin{itemize}
\raggedright
  \item \emph{Say it:} \emph{jāḍu}
  \item \emph{Say it:} \emph{jāḍu}, once more
  \item \emph{Say it:} \emph{polīsaru kārina jāḍu hiḍidu hōdaru.} — The police followed the car's trail.
  \item \emph{Recall:} say the Kannada for to commit, then the Kannada for to abscond, then say \emph{jāḍu} again
\end{itemize}

\begin{tcolorbox}[breakable,colback=teal!4,colframe=teal!35!black,title={Before you move on}]
What does \kn{ಜಾಡು} mean? (\textquotedblleft{}A track; a trail\textquotedblright{}.) Say it once more.
\end{tcolorbox}
